% =============================================================================
% main.tex  --  thesis body, \input from the scrbook template
%
% Every number in a table below was measured in this project. Experiments that
% have not been run yet are marked \TODO{...} so nothing is fabricated. Search
% for "TODO" before submission and fill each one from the corresponding log.
% =============================================================================

% --- \TODO, \mm and the packages used below are defined in thesis.tex ------

% =============================================================================
\chapter{Introduction}
\label{ch:intro}
% =============================================================================

\section{Problem}

Human Mesh Recovery (HMR) is the task of estimating the full 3D surface of the
people visible in an image or video. Rather than regressing thousands of
vertices directly, HMR methods predict the low-dimensional parameters of a
\emph{parametric body model}: a differentiable function that maps a pose vector
and a shape vector to a watertight mesh. The parametric model supplies a strong
anatomical prior, makes the output animatable, and lets a network learn from a
compact target.

Almost every HMR system published in the last decade targets the SMPL family of
body models~\cite{loper2015smpl,romero2017mano,pavlakos2019smplx}. Two
properties of that family motivate this thesis. First, SMPL is released under a
research-only license that forbids commercial use. Second, its shape space is
learned by principal component analysis over 3D scans of adults, so it does not
faithfully represent children, the elderly, or the extremes of body type, and
its ten shape coefficients carry no semantic meaning.

Anny~\cite{bregier2025anny} is a recently released body model that addresses
both points. It is published under the Apache~2.0 license, and its shape is
controlled by \emph{phenotype} parameters, e.g.\ \texttt{age}, \texttt{gender},
\texttt{height}, \texttt{weight}, \texttt{proportions}, that are interpretable
and calibrated against WHO population statistics. Anny is accompanied by
Anny-One, a large synthetic dataset of photorealistic renderings with exact
ground truth. No published multi-person HMR system predicts Anny parameters.

The problem addressed here is therefore: \emph{adapt a strong multi-person HMR
architecture so that it outputs Anny parameters, train it on Anny-One, and
reduce its size substantially while retaining accuracy}. The research question
posed by the supervising group is whether an accurate and efficient multi-person
HMR system for Anny can be trained, with a model size significantly smaller than
a strong baseline.

\section{Motivation}

Multi-HMR~\cite{baradel2024multihmr} is a single-shot, multi-person, whole-body
HMR network: one forward pass of a Vision Transformer (ViT) detects every person
in the image and regresses their SMPL-X parameters and 3D location. It is an
attractive base because detection and mesh regression share one backbone, the
architecture is body-model agnostic in principle (the body model appears only
in the final layer and in the loss), and public checkpoints with several
backbone sizes exist. Its best configuration, however, uses a ViT-L backbone
of roughly 300\,M parameters, which is far from the efficiency target.

Adapting it to Anny is not a one-line change. Anny has 163 bones where SMPL-X
has 55, expresses pose as $4{\times}4$ transforms rather than axis-angle vectors,
and has a shape space with no correspondence to SMPL-X's PCA coefficients.
Meeting the size target additionally requires a compression strategy that does
not simply discard the accuracy of the large model.

\section{Contributions}

\begin{enumerate}
  \item A working Multi-HMR variant that predicts Anny parameters (163 bone
        rotations in a 6D representation plus 11 phenotypes), trained on
        Anny-One, reaching 74.3\,mm PVE on a held-out split from a 145\,mm
        starting point via a staged backbone-transfer schedule
        (Chapter~\ref{ch:experiments}).
  \item A knowledge-distillation framework that compresses the ViT-L model
        8.6$\times$ into a ViT-S student. Feature-level distillation alone
        left the student at 193\,mm PVE; transferring the teacher's prediction
        heads through a learned projector that feeds them teacher-width tokens
        brings it to 127\,mm, with each intermediate step measured
        (Section~\ref{sec:kd}, Section~\ref{sec:kd_results}).
  \item Two measured defects and their fixes: a shape head that regressed
        dataset averages because of loss-weight imbalance, and a neck
        elongation artifact whose cause was isolated by eliminating four
        alternative hypotheses (Chapter~\ref{ch:ablation}).
  \item An optimisation-based converter from SMPL ground truth to Anny
        parameters whose joint correspondence is verified by rest-pose
        geometry rather than by name, and which closes the remaining
        definition gap between the two rigs with \emph{virtual joints}. The
        error floor fell from 34\,mm to 13--20\,mm, and the whole of 3DPW
        (87 person--sequences) was converted into Anny-space labels for mixed
        synthetic--real training (Section~\ref{sec:smpl2anny},
        Section~\ref{sec:fit_results}).
  \item A reproducible experimental setup in which every intervention is a
        command-line flag, together with a documented set of pipeline bugs
        (silent weight loading, unit mismatches, per-image intrinsics) whose
        symptoms are easily mistaken for model failures
        (Section~\ref{sec:pipeline}).
\end{enumerate}

\section{Outline}

Chapter~\ref{ch:related} reviews parametric body models, HMR architectures and
knowledge distillation, and states the research gap. Chapter~\ref{ch:method}
gives the formal problem statement, the network, its inputs and outputs, and the
loss functions. Chapter~\ref{ch:experiments} describes datasets, metrics,
baselines and the main results. Chapter~\ref{ch:ablation} presents ablations
and the two root-cause investigations. Chapter~\ref{ch:discussion} discusses
what did and did not improve and why, and Chapter~\ref{ch:conclusion} concludes.

% =============================================================================
\chapter{Related Work}
\label{ch:related}
% =============================================================================

\section{Parametric body models}

A parametric body model is a differentiable function
$M(\boldsymbol{\theta},\boldsymbol{\beta}) \rightarrow \mathbf{V}\in\mathbb{R}^{N\times 3}$
that maps a pose vector $\boldsymbol{\theta}$ and a shape vector
$\boldsymbol{\beta}$ to $N$ mesh vertices. Pose is applied by linear blend
skinning (LBS): every vertex is a weighted combination of the transforms of a
few nearby bones, so the mesh deforms smoothly as the skeleton moves.

\paragraph{SMPL~\cite{loper2015smpl}.} The pioneering model of this family
(2015) has 6,890 vertices, 24 joints, and a shape space of 10 PCA coefficients
learned from thousands of registered scans. Its parameters are a global
orientation (3), body pose (69 axis-angle values for 23 joints), translation
(3), and betas (10). SMPL-H~\cite{romero2017mano} adds 15 articulated joints per
hand, and SMPL-X~\cite{pavlakos2019smplx} adds an expressive face (jaw, eyes and
10 expression coefficients), giving 55 joints in total. Multi-HMR predicts
SMPL-X.

\paragraph{Anny~\cite{bregier2025anny}.} Anny (2025) departs from the
scan-based design. Its mesh (13,718 vertices, 13,710 quad faces) and rig of
163 bones come from the open MakeHuman project, and its shape is a set of
artist-authored blendshapes controlled by \emph{phenotypes}. The subset used
in this work is \texttt{gender}, \texttt{age}, \texttt{muscle}, \texttt{weight},
\texttt{height}, \texttt{proportions}, \texttt{cupsize}, \texttt{firmness} and
three ethnicity blends (\texttt{african}, \texttt{asian}, \texttt{caucasian}),
each in $[0,1]$. Pose is a $4{\times}4$ transform per bone. Because the model
is built from open assets rather than scans it is free of biometric-data
constraints and is released under Apache~2.0. The authors show that HMR
models trained with Anny-One match those trained with scan-based models.

Two properties of Anny's rig matter later in this thesis. Of the 163 bones,
roughly 88 are \emph{helper} or \emph{deform} bones (breast, spine and neck
subdivisions, face micro-bones) whose rotation cannot be inferred from an
image; and the neck is modelled by three stacked bones
(\texttt{neck01}, \texttt{neck02}, \texttt{neck03}) plus \texttt{head}.

\section{Human mesh recovery}

Early HMR methods regressed SMPL parameters for a single cropped person, most
influentially HMR~\cite{kanazawa2018hmr} with an iterative regression head.
HMR2.0~\cite{goel2023humans4d} replaced the convolutional backbone with a
Vision Transformer and showed that a large ViT backbone trained on a broad
mixture of data is a strong, simple recipe, and extended it to video by tracking
people in 3D.

Multi-person HMR adds detection. Two-stage methods run a person detector and
then a single-person HMR on each crop, which is slow and loses inter-person
context. Single-shot methods such as ROMP~\cite{sun2021romp} and
BEV~\cite{sun2022bev} predict a person-centre heatmap and per-location
parameters in one pass.

\paragraph{Multi-HMR~\cite{baradel2024multihmr}.} This is the base of the
present work. A ViT backbone (DINOv2~\cite{oquab2023dinov2} initialisation)
produces one token per $14{\times}14$ image patch. A small MLP scores each
patch as ``a person's primary keypoint is here'', and detected tokens become
queries to the \emph{Human Perception Head} (HPH), a stack of cross-attention
blocks in which every query attends to the full image and to the other queries.
The HPH outputs SMPL-X parameters, a 2D offset, and a log-depth from which the
3D translation is recovered through the camera intrinsics. Multi-HMR is trained
on synthetic data (BEDLAM~\cite{black2023bedlam}, AGORA~\cite{patel2021agora}
and the authors' close-up dataset CUFFS) and reports, with a ViT-B backbone and
HPH head, 73.5\,mm MPJPE on 3DPW; a ViT-S backbone at $448{\times}448$ runs at
roughly 30\,fps and a ViT-L at $896{\times}896$ at roughly 5\,fps on a V100.

\section{Knowledge distillation}

Knowledge distillation~\cite{hinton2015distilling} trains a small
\emph{student} network to reproduce the behaviour of a large, already-trained
\emph{teacher}. The original formulation matches the teacher's softened
output probabilities with a Kullback--Leibler (KL) divergence at temperature
$T$, scaling the gradient by $T^2$ so that its magnitude is independent of
$T$. \emph{Feature} distillation~\cite{romero2015fitnets} instead matches
intermediate representations, which is the form used here: the student
backbone is asked to reproduce the teacher backbone's patch tokens, and a
learned linear projector bridges the different channel widths. Distillation
is the standard way to obtain a small model that retains most of a large
model's accuracy, and it is directly applicable to Multi-HMR because the
backbone is a separable module in front of the heads.

\section{Research gap}

Three gaps are addressed.
\begin{enumerate}
  \item \textbf{No HMR network predicts Anny.} All multi-person HMR systems
        output SMPL-family parameters. Using Anny requires redesigning the
        output layer for 163 rotations and semantic phenotypes, converting the
        ground truth, and re-examining the loss, since Anny's helper bones and
        phenotype space have no SMPL-X analogue.
  \item \textbf{Model size.} The best Multi-HMR configuration uses a ViT-L
        backbone of $\sim$300\,M parameters. The thesis target is a
        25--50\,\% or larger reduction at similar accuracy, which the paper's
        own ablations show a from-scratch ViT-S does not achieve (80.1 vs.\
        73.5\,mm MPJPE on 3DPW for ViT-S vs.\ ViT-B). Distillation is the
        natural tool and has not been applied to this architecture.
  \item \textbf{Real-image evaluation in Anny space.} Anny-One is synthetic.
        Real-world benchmarks such as 3DPW~\cite{vonmarcard2018_3dpw} carry
        SMPL ground truth, so evaluating or training an Anny model on them
        requires converting SMPL parameters into Anny's space, which no
        existing tool does.
\end{enumerate}

% =============================================================================
\chapter{Method}
\label{ch:method}
% =============================================================================

\section{Problem statement}

Given an RGB image $\mathbf{I}\in\mathbb{R}^{H\times W\times 3}$ with (optionally
known) intrinsics $\mathbf{K}\in\mathbb{R}^{3\times3}$, containing an unknown
number $n$ of people, predict for each person $i\in\{1,\dots,n\}$:
\begin{align}
  \mathbf{R}_i &\in SO(3)^{J},\quad J=163 \quad\text{(per-bone rotations)}\\
  \boldsymbol{\phi}_i &\in [0,1]^{P},\quad P=11 \quad\text{(phenotypes)}\\
  \mathbf{t}_i &\in \mathbb{R}^3 \quad\text{(3D translation in camera coordinates)}
\end{align}
such that the Anny mesh
$\mathbf{V}_i = M_{\text{Anny}}(\mathbf{R}_i,\boldsymbol{\phi}_i) + \mathbf{t}_i$
matches the person's true body surface. A secondary objective is to do this with
a network whose parameter count is a small fraction of the ViT-L baseline.

\section{Network architecture}
\label{sec:arch}

The architecture (Figure~\ref{fig:arch}) follows Multi-HMR; the modifications
are confined to the parameter head and the body-model layer.

\paragraph{Backbone.} A DINOv2 ViT with patch size 14 maps a
$672{\times}672$ input to a $48{\times}48$ grid of $C$-dimensional tokens
$\mathbf{z}\in\mathbb{R}^{2304\times C}$ ($C{=}1024$ for ViT-L, 768 for ViT-B,
384 for ViT-S).

\paragraph{Detection.} A two-layer MLP scores every token; after non-maximum
suppression, tokens above a threshold are treated as detected people. A second
MLP predicts a 2D sub-patch offset so the person's primary keypoint (the head)
is localised more finely than the patch grid.

\paragraph{Human Perception Head.} Detected tokens become queries to a
cross-attention transformer (dimension 512, 2 blocks, 16 heads). Each query
attends to all image tokens and, through self-attention, to the other people.
The output token $\mathbf{y}_i$ feeds three MLPs.

\paragraph{Anny parameter head.} This is the main architectural change.
\begin{itemize}
  \item \emph{Pose.} Rotations are regressed in the continuous 6D
        representation of Zhou et al.~\cite{zhou2019continuity}: the MLP
        outputs $\mathbf{r}_i\in\mathbb{R}^{163\times 6}$ as a residual on a
        fixed mean-pose buffer, and each 6-vector is mapped to a rotation
        matrix by Gram--Schmidt orthonormalisation. This replaces SMPL-X's
        $55{\times}3$ axis-angle output, avoiding the discontinuities of
        axis-angle regression.
  \item \emph{Helper-bone masking.} A fixed binary mask
        $\mathbf{m}\in\{0,1\}^{163}$ marks 75 bones as predictable and forces
        the remaining 88 helper bones to identity,
        $\mathbf{R}_{i,j}\leftarrow m_j\mathbf{R}_{i,j}+(1-m_j)\mathbf{I}$.
        The same mask is applied to the ground truth so those bones contribute
        zero loss.
  \item \emph{Shape.} An MLP outputs $P{=}11$ logits passed through a sigmoid to
        give phenotypes in $[0,1]$, replacing the SMPL-X beta/expression heads.
  \item \emph{Depth.} As in Multi-HMR, a log-depth $d_i$ normalised by focal
        length is predicted and unprojected with $\mathbf{K}$ to obtain
        $\mathbf{t}_i$.
\end{itemize}

\paragraph{Body-model layer.} The Anny model is evaluated inside the network,
in PyTorch, to produce vertices $\mathbf{V}_i$ and joint positions
$\mathbf{J}_i$ from $(\mathbf{R}_i,\boldsymbol{\phi}_i)$, so that 3D and 2D
geometric losses can be back-propagated through it.

\begin{figure}[t]
  \centering
  \fbox{\parbox{0.9\linewidth}{\centering\vspace{2cm}
  \TODO{insert architecture figure: image $\rightarrow$ ViT backbone
  $\rightarrow$ detection MLP + HPH $\rightarrow$ Anny head (6D pose, phenotypes,
  depth) $\rightarrow$ Anny body model $\rightarrow$ mesh; teacher backbone in
  a dashed box feeding the KD loss (Section~\ref{sec:kd})}\vspace{2cm}}}
  \caption{Overview of the Anny-adapted Multi-HMR architecture.}
  \label{fig:arch}
\end{figure}

\section{Inputs and outputs}

\begin{table}[h]
\centering\small
\begin{tabular}{lll}
\toprule
 & Tensor & Shape \\
\midrule
Input  & RGB image & $3{\times}672{\times}672$ \\
       & intrinsics $\mathbf{K}$ (optional) & $3{\times}3$ \\
\midrule
Output & detection scores & $48{\times}48$ \\
(per person) & 2D offset & $2$ \\
       & bone rotations $\mathbf{R}$ & $163{\times}3{\times}3$ \\
       & phenotypes $\boldsymbol{\phi}$ & $11$ \\
       & log-depth $d$, translation $\mathbf{t}$ & $1$, $3$ \\
       & mesh vertices $\mathbf{V}$ & $13718{\times}3$ \\
       & joints $\mathbf{J}$ & $163{\times}3$ \\
\bottomrule
\end{tabular}
\caption{Inputs and outputs of the network.}
\end{table}

\section{Loss functions}
\label{sec:loss}

Ground truth is prepared per person: Anny $4{\times}4$ transforms are converted
to rotation matrices, and the Anny model is run on the ground-truth parameters
to obtain target vertices and joints. The total loss is
\begin{equation}
\begin{split}
\mathcal{L} = {}& \alpha_{\text{bce}}\mathcal{L}_{\text{det}}
  + \alpha_{\text{off}}\mathcal{L}_{\text{off}}
  + \alpha_{\text{rot}}\mathcal{L}_{\text{rot}}
  + \alpha_{\text{shape}}\mathcal{L}_{\text{shape}}
  + \alpha_{\text{dist}}\mathcal{L}_{\text{dist}}
  + \alpha_{\text{t}}\mathcal{L}_{\text{t}} \\
  &+ \alpha_{\text{j3d}}\mathcal{L}_{\text{j3d}}
  + \alpha_{\text{v3d}}\mathcal{L}_{\text{v3d}}
  + \mathbb{1}[e \ge e_{2\text{D}}]\left(\alpha_{\text{j2d}}\mathcal{L}_{\text{j2d}}
  + \alpha_{\text{v2d}}\mathcal{L}_{\text{v2d}}\right)
  + \lambda_{\text{kd}}\mathcal{L}_{\text{kd}}
\end{split}
\label{eq:loss}
\end{equation}
with
\begin{align}
\mathcal{L}_{\text{det}} &= \text{focal loss on the person-centre heatmap (CenterNet form)}\\
\mathcal{L}_{\text{rot}} &= \sum_{j} w_j \left\| \mathbf{R}_j - \hat{\mathbf{R}}_j \right\|_1,
  \qquad \textstyle\frac{1}{J}\sum_j w_j = 1\\
\mathcal{L}_{\text{shape}} &= \left\| \boldsymbol{\phi} - \hat{\boldsymbol{\phi}} \right\|_1\\
\mathcal{L}_{\text{j3d}} &= \frac{1}{J}\sum_j w_j\left\| (\mathbf{J}_j - \mathbf{J}_{\text{pelvis}}) - (\hat{\mathbf{J}}_j - \hat{\mathbf{J}}_{\text{pelvis}}) \right\|_1\\
\mathcal{L}_{\text{v3d}} &= \frac{1}{N}\sum_v \left\| (\mathbf{V}_v - \mathbf{J}_{\text{pelvis}}) - (\hat{\mathbf{V}}_v - \hat{\mathbf{J}}_{\text{pelvis}}) \right\|_1
\end{align}
and $\mathcal{L}_{\text{j2d}},\mathcal{L}_{\text{v2d}}$ the $\ell_1$ reprojection
errors of joints and vertices, enabled only after epoch $e_{2\text{D}}$ because
early predictions are too poor for reprojection gradients to be informative.
$\mathcal{L}_{\text{off}},\mathcal{L}_{\text{dist}},\mathcal{L}_{\text{t}}$ are
$\ell_1$ losses on the offset, log-depth and translation. Hats denote ground
truth. The per-joint weights $w_j$ are $1$ by default; Section~\ref{sec:neck}
uses them to up-weight neck and hand bones.

The Multi-HMR defaults are
$\alpha_{\text{bce}}{=}10,\ \alpha_{\text{rot}}{=}0.1,\ \alpha_{\text{shape}}{=}1,\
\alpha_{\text{j3d}}{=}\alpha_{\text{v3d}}{=}100$. The 100:1 ratio between the
geometric terms and the shape term is the subject of
Section~\ref{sec:shape_ablation}, where $\alpha_{\text{shape}}{=}20$ is shown to
be necessary for Anny.

\section{Knowledge distillation}
\label{sec:kd}

To meet the size target, a frozen ViT-L \emph{teacher} (the best trained model)
supervises a ViT-S \emph{student} at the level of backbone patch tokens. The
loss is inserted immediately after patch decoding and before any prediction
head, so the student is asked to reproduce the teacher's \emph{representation}
rather than only its outputs. With teacher tokens
$\mathbf{z}^{t}\in\mathbb{R}^{N\times C_t}$ and student tokens
$\mathbf{z}^{s}\in\mathbb{R}^{N\times C_s}$, a learned linear projector
$\mathbf{P}\in\mathbb{R}^{C_s\times C_t}$ bridges the width mismatch
($C_s{=}384$, $C_t{=}1024$), and
\begin{equation}
\mathcal{L}_{\text{kd}} = T^2\cdot \frac{1}{N}\sum_{n=1}^{N}
  \mathrm{KL}\!\left(\ \sigma\!\left(\tfrac{\mathbf{z}^{t}_n}{T}\right)
  \ \Big\|\ \sigma\!\left(\tfrac{\mathbf{z}^{s}_n\mathbf{P}}{T}\right)\right),
\label{eq:kd}
\end{equation}
where $\sigma$ is the softmax over the channel dimension and $T$ the
temperature~\cite{hinton2015distilling}. The $T^2$ factor keeps the gradient
magnitude comparable across temperatures. The teacher receives no gradient; the
projector is trained with the student and discarded at inference, so the
deployed model is the student alone. $\mathcal{L}_{\text{kd}}$ is added to
Equation~\ref{eq:loss} with weight $\lambda_{\text{kd}}$; the student is trained
on the full task loss at the same time.

\paragraph{Head transfer.} The first student (Section~\ref{sec:kd_results})
showed that feature-level distillation alone leaves the prediction heads, which
are trained from random initialisation, far behind the teacher's. Because
teacher and student instantiate the same \texttt{Model} class, every head
tensor has the same name in both; the student is therefore initialised by
copying each non-backbone tensor whose name \emph{and shape} match the
teacher's (\texttt{--init\_heads\_from\_teacher}). Of 67 non-backbone tensors,
54 transfer. The remaining 13 are exactly those whose shape depends on the
backbone width $C$ ($1024$ vs $384$): the HPH token embedding, the
cross-attention key/value projections in both HPH blocks, and the detection
and offset MLPs. These \emph{entry points} of the head remain randomly
initialised.

\paragraph{Output-level distillation.} In addition to matching backbone
tokens, the student's predictions are pulled toward the teacher's on the same
image. Both models are run on the ground-truth query locations during
training, so person $i$ of the student corresponds to person $i$ of the
teacher without any matching step:
\begin{equation}
\mathcal{L}_{\text{kd-out}} = \alpha_{\text{rot}}\left\|\mathbf{R}^{s}-\mathbf{R}^{t}\right\|_1
 + \alpha_{\text{shape}}\left\|\boldsymbol{\phi}^{s}-\boldsymbol{\phi}^{t}\right\|_1
 + \alpha_{\text{dist}}\left|d^{s}-d^{t}\right| ,
\label{eq:kdout}
\end{equation}
weighted like the task terms and added with $\lambda_{\text{kd-out}}$
(\texttt{--lambda\_kd\_out}). Teacher outputs are detached.

\paragraph{Warm-up freeze.} With trained head blocks reading through random
entry projections, the head produces confident but wrong outputs from step
one, and the resulting gradients degrade the transferred weights before the
projections have adapted. The transferred tensors are therefore frozen for the
first $E_w$ epochs (\texttt{--freeze\_heads\_epochs}), during which only the 13
random projections, the backbone and the KD projector train, with
Equation~\ref{eq:kdout} supplying a direct target; all parameters are then
released. The 2D reprojection losses are enabled at the same epoch.


\paragraph{Projector-fed heads.} The warm-up freeze protects the transferred
weights but cannot help the 13 random entry projections, which must learn, from
384-dimensional student tokens, to reproduce what the teacher's own projections
produced from 1024-dimensional tokens. Meanwhile the KD projector $\mathbf{P}$
of Equation~\ref{eq:kd} learns exactly this mapping, yet lives inside the loss
and is discarded. The final design therefore moves the projector \emph{in front
of the heads}: the student's tokens are mapped to the teacher width once,
$\mathbf{z}^{s}\mathbf{P}$, every head is built at that width
(\texttt{--head\_dim 1024}), and the feature loss of Equation~\ref{eq:kd}
compares the projected tokens with the teacher's directly, without a
loss-side projector. All 67 head tensors then transfer, the only new parameters
are the $384{\times}1024$ projector ($0.4$\,M), and the heads receive inputs of
the kind they were trained on from the first step, while the feature loss
trains the projector toward precisely that input distribution. The deployed
student is backbone, projector and heads: 37\,M parameters, 8.6$\times$ smaller
than the teacher.

\paragraph{Throughput.} The frozen teacher dominates each distillation
iteration: a ViT-L at $672^2$ processes 2304 tokens through 24 blocks,
$\approx$1.9\,TFLOP per image, and with mixed precision disabled (it produced
NaNs with the Anny layer) these matmuls ran in plain fp32. Two changes that
leave the student's own arithmetic untouched cut the iteration time from
1.24\,s to 0.47\,s on an A100: enabling TensorFloat-32 matmuls (fp32 range,
10-bit mantissa inputs, fp32 accumulation), and running the teacher's backbone
under bfloat16 autocast inside \texttt{no\_grad}, casting its tokens back to
fp32 before its heads and the loss. The first evaluation after the change
reproduced the previous one to the decimal (155.0/107.3\,mm), confirming that
only speed changed.

\section{SMPL-to-Anny conversion for real data}
\label{sec:smpl2anny}

Real-world datasets such as 3DPW carry SMPL ground truth. To train and
evaluate in Anny space, every annotated person is converted by optimisation,
following the autograd strategy of NVlabs' \texttt{smpl2soma} and the
joint-fitting objective of \texttt{joints2smpl}: Anny's 163 6D rotations
$\mathbf{r}$, phenotype logits $\boldsymbol{\ell}$ and a translation
$\mathbf{t}$ are found so that Anny's joints land on the 24 SMPL joints
$\hat{\mathbf{J}}^{S}$ of the frame,
\begin{equation}
\min_{\mathbf{r},\boldsymbol{\ell},\mathbf{t}}
\ \frac{1}{24}\sum_{k=1}^{24}
\left\| \mathbf{p}_{\pi(k)}\!\left(\mathbf{r},\sigma(\boldsymbol{\ell})\right)
 + \mathbf{t} - \hat{\mathbf{J}}^{S}_{k}\right\|_2^2
 + \lambda_{r}\sum_{b\neq\text{root}}\left\|\mathbf{R}_b-\mathbf{I}\right\|_F^2
 + \lambda_{\phi}\sum_{j\in\text{free}}\left(\sigma(\ell_j)-\tfrac12\right)^2 ,
\label{eq:fit}
\end{equation}
where $\pi$ is the bone correspondence and $\mathbf{p}_{b}$ the position of the
point attached to bone $b$ (defined below). Optimisation is staged (root
orientation and translation, then all bone rotations, then phenotypes), with
Adam on the 6D representation. The two regularisers are weak
($\lambda_r{=}10^{-4}$, $\lambda_\phi{=}10^{-3}$): the first settles the null
space of the problem (twist about a bone's own axis moves no joint but
corkscrews the mesh), the second stops a single residual joint gap from
dragging a phenotype to an extreme. Only the root and the bones that carry a
mapped joint may rotate; the other $\sim$130 bones (twist helpers, fingers,
toes, face) have no target and would otherwise drift.

\paragraph{Targets.} 3DPW's own per-frame \texttt{jointPositions} are used
as $\hat{\mathbf{J}}^{S}$ rather than a neutral-model SMPL forward pass:
3DPW's shape coefficients belong to the gendered SMPL models, and evaluating
them in the neutral model moved the joints by 37\,mm on average (93\,mm at the
feet) on the first frame checked.

\paragraph{Correspondence verified by geometry.} The first version of the
converter chose Anny bones by name. Two of those choices were wrong in a way
that no name reveals: Anny's MakeHuman rig numbers the spine \emph{top-down}
(\texttt{spine01} sits just below the neck, \texttt{spine05} on the pelvis)
while SMPL numbers it bottom-up, so SMPL's lumbar \texttt{spine1} had been
mapped to a bone 25\,cm above it; and SMPL's hip joints had been mapped to
\texttt{pelvis.L/R}, short bones whose heads lie 4\,cm from the hip joint,
instead of to \texttt{upperleg01.L/R}. The converter now evaluates both
models in their rest poses, prints the two skeletons side by side, and
chooses the spine and hip bones by their height along the hip-midpoint--neck
axis and their distance from the pelvis.

\paragraph{Virtual joints.} Even with the right bones, the two rigs define
their torso joints at different anatomical points: in the rest pose SMPL's
pelvis joint sits 9\,cm above its hip joints whereas Anny's root sits at hip
height and 7\,cm in front of them, SMPL's \texttt{spine3} lies between two Anny
bones 6.5\,cm from either, and its collar joints are 6\,cm further lateral.
Forcing bone \emph{heads} onto such targets makes the optimiser compromise
with the pelvis and the shape. Instead, for each torso joint the rest-pose
offset from its Anny bone head to the SMPL joint is measured once (both
skeletons aligned at the hip midpoint, with the spine as the vertical axis
and scaled to equal hip--neck length) and expressed in the bone's local frame
$\boldsymbol{\delta}_b$; the fit then moves the point
$\mathbf{p}_b = \mathbf{h}_b + \mathbf{R}_b\boldsymbol{\delta}_b$ onto the SMPL
joint, where $\mathbf{h}_b$ and $\mathbf{R}_b$ are the posed bone head and
frame. The offsets are 11.6\,cm for the pelvis, 4.3\,cm for the hips, 3--7\,cm
along the spine and 5\,cm for the collars; limb joints use the bone head.

\paragraph{Shape.} Phenotypes are not \emph{mapped} from SMPL's PCA
coefficients, which have no semantic counterpart. Nor can joint positions
determine most of them: weight, muscle, cup size, firmness and ethnicity move
no joint, and adult age barely does. Left free, the optimiser used them as
slack and produced the same body for every actor (all four test-scene
subjects came out female, 0.34 in height, and all adults at exactly 0.67 in
age). The final converter therefore fixes gender from 3DPW's annotation, fixes
age at 0.55 (the Anny-One adult mean; the 3DPW actors are young adults),
fits only height and proportions, and estimates one body per person and
sequence as the per-phenotype median over all frames before refitting every
frame with that body fixed and pose free.

\paragraph{Whole-dataset pipeline.} 3DPW is video, so each frame is
warm-started from the previous one, which cuts the per-frame cost to
$\sim$8\,s (every third frame is converted, 10\,fps). The 87
person--sequences of the train, validation and test splits are converted as
one Slurm array job; for each frame the posed mesh is projected with 3DPW's
camera to obtain a bounding box and the fraction of the body inside the
image. A final pass merges the people of each frame into one label file per
image in the Anny-One format (camera intrinsics, and per person the
$163{\times}4{\times}4$ pose with the root carried into the camera frame as
$\mathbf{R}_{\text{cam}}\mathbf{Q}_{\text{root}}$ and
$\mathbf{R}_{\text{cam}}(\mathbf{t}+\mathbf{t}_{\text{3DPW}})+\mathbf{t}_{\text{cam}}$,
the 11 phenotypes and the box), dropping people whose fit error exceeds
80\,mm or who are mostly outside the frame. The camera convention is the one
Multi-HMR's own 3DPW loader uses and was checked independently by projecting
the converted joints onto 3DPW's 2D detections (8.9\,px mean). A dataset class
returns these labels in the exact sample format of Anny-One, so the training
code, collate function and evaluation run on them unchanged.

\paragraph{Mixed synthetic--real training.} With the labels in place, the
student is trained on the union of Anny-One and the converted 3DPW-train
split, drawn by a weighted sampler so that on average a fraction $\rho$ of
every batch is real (\texttt{--mix\_3dpw\_ratio}, $\rho{=}0.3$), and evaluated
every few epochs on both the Anny-One holdout and the converted
3DPW-validation split. The 3DPW test split is never used for training.

\section{Training and transfer schedule}
\label{sec:schedule}

Training of the teacher begins from the public Multi-HMR ViT-L checkpoint
\texttt{multiHMR\_672\_L\_anny}; the student backbone is analogously initialised
from the authors' \texttt{multiHMR\_672\_S}, a ViT-S already trained for mesh
recovery, rather than from ImageNet-only DINOv2. Because the Anny
head has a different architecture from the SMPL-X head, only the backbone is
transferred (a key-name remapping is required between the original and the
training model class). A staged schedule is used:
\begin{enumerate}
  \item \emph{Frozen backbone:} only the heads train (learning rate $3{\cdot}10^{-5}$).
  \item \emph{Partial fine-tuning:} the last 4 of 24 ViT blocks are unfrozen
        ($10^{-5}$), with helper-bone masking enabled.
  \item \emph{Loss re-weighting:} $\alpha_{\text{shape}}$ raised to 20.
\end{enumerate}
Every continuation uses a learning rate one order of magnitude below the
cold-start value, with step decay (halving every 30--60 epochs); a constant
cold-start rate on a partially trained model produced hundreds of NaN gradient
events per epoch. All of these choices are command-line flags
(\texttt{--freeze\_backbone}, \texttt{--unfreeze\_last\_n\_blocks},
\texttt{--mask\_helper\_joints\_train}, \texttt{--alpha\_shape},
\texttt{--boost\_neck\_weight}, \texttt{--distill\_teacher\_ckpt},
\texttt{--lambda\_kd}, \texttt{--kd\_temperature}) so that any configuration can
be reproduced without editing code.

% =============================================================================
\chapter{Experiments}
\label{ch:experiments}
% =============================================================================

\section{Datasets}

\paragraph{Anny-One~\cite{bregier2025anny} (training and in-domain test).}
579,960 photorealistic $672{\times}672$ renderings of Anny bodies in indoor and
outdoor scenes, each with intrinsics $\mathbf{K}$ and, for every visible person,
the exact Anny pose ($163{\times}4{\times}4$) and phenotypes. The last 100
samples are held out and never trained on; they serve as the in-domain test set
(\texttt{--test\_anny\_n 100}). Per-image field of view varies widely
(51$^\circ$--110$^\circ$ across a sample of five images), so the true
$\mathbf{K}$ must be used at evaluation.

\paragraph{3DPW~\cite{vonmarcard2018_3dpw} (real-world training and evaluation).}
60 outdoor and indoor video sequences ($1080{\times}1920$, 30\,fps) with SMPL
ground truth from IMUs and a hand-held camera, in the licensed standard
release: 24 training, 12 validation and 24 test sequences, 87 annotated
person--sequences in total. All of them are converted to Anny space with the
method of Section~\ref{sec:smpl2anny} at every third frame. The train split is
used for mixed training, the validation split for model selection and the
test split is reserved for reporting. Two evaluation protocols exist and are
reported separately: (a) the Multi-HMR protocol against the original SMPL
meshes, comparable to published numbers; (b) Anny-space metrics against the
converted ground truth, on which the conversion residual (13--20\,mm,
Section~\ref{sec:fit_results}) is a floor.

\section{Evaluation metrics}
\label{sec:metrics}

For a person with predicted joints $\mathbf{J}\in\mathbb{R}^{J\times3}$ and
vertices $\mathbf{V}\in\mathbb{R}^{N\times3}$, and ground truth
$\hat{\mathbf{J}},\hat{\mathbf{V}}$, both pelvis-centred:
\begin{align}
\text{MPJPE} &= \frac{1}{J}\sum_{j=1}^{J}\left\|\mathbf{J}_j-\hat{\mathbf{J}}_j\right\|_2\\
\text{PVE}   &= \frac{1}{N}\sum_{v=1}^{N}\left\|\mathbf{V}_v-\hat{\mathbf{V}}_v\right\|_2\\
\text{PA-MPJPE} &= \frac{1}{J}\sum_{j=1}^{J}\left\|s\mathbf{R}\mathbf{J}_j+\mathbf{t}-\hat{\mathbf{J}}_j\right\|_2,
\quad (s,\mathbf{R},\mathbf{t}) = \arg\min \sum_j \left\|s\mathbf{R}\mathbf{J}_j+\mathbf{t}-\hat{\mathbf{J}}_j\right\|_2^2
\end{align}
and PA-PVE analogously on vertices. The Procrustes-aligned (PA) variants remove
global rotation, scale and translation and therefore measure the quality of the
body reconstruction alone; the unaligned variants additionally penalise
orientation error. All are reported in millimetres. Because these metrics are
computed only over predictions matched to a ground-truth person, detection
precision, recall and F1 are reported alongside them: a model that detects few
people can post a good PVE while missing most of them.

\section{Baselines}

Baseline numbers are taken from the Multi-HMR paper~\cite{baradel2024multihmr}
(Table~\ref{tab:baselines}). They are SMPL-X predictions on 3DPW and are not
directly comparable to Anny-space numbers, but they establish the accuracy
gap between backbone sizes that distillation must close.

\begin{table}[h]
\centering\small
\begin{tabular}{llrr}
\toprule
Backbone & Head & Params (backbone) & 3DPW MPJPE $\downarrow$ \\
\midrule
ViT-S & HPH & 21\,M  & 80.1 \\
ViT-B & HPH & 86\,M  & 73.5 \\
ViT-L & HPH & 304\,M & \TODO{from paper Fig.~4 / suppl.} \\
\bottomrule
\end{tabular}
\caption{Multi-HMR baselines (SMPL-X, trained on AGORA+BEDLAM) reproduced from
\cite{baradel2024multihmr}.}
\label{tab:baselines}
\end{table}

\section{Implementation details}

Input $672{\times}672$, AdamW, 1000 iterations per epoch, evaluation every 5
epochs on the holdout, mixed precision disabled for trainable models (it
produced NaNs with the Anny layer). Teacher runs use batch size 4; student runs
batch 8 (the frozen teacher needs no activation memory), with TF32 matmuls and
a bfloat16 teacher backbone (Section~\ref{sec:kd}), at 0.47\,s per iteration
on one NVIDIA A100 of the DFKI Pegasus cluster. Model sizes: ViT-L 318.4\,M
parameters; ViT-S student with 384-wide heads 32.8\,M, with projector-fed
1024-wide heads 37\,M (8.6$\times$ smaller).

\section{Main results: Anny-One}
\label{sec:main_results}

Table~\ref{tab:main} traces the held-out PVE across the transfer schedule.
The frozen-backbone stage reaches 145\,mm; unfreezing the last four blocks and
masking helper bones halves this to 73.6\,mm, after which further loss
re-weighting changes PVE only marginally but substantially improves shape
estimation (Section~\ref{sec:shape_ablation}).

\begin{table}[h]
\centering\small
\begin{tabular}{llrrrrr}
\toprule
Run & Change vs.\ previous & PVE & PA-PVE & MPJPE & PA-MPJPE & Recall \\
\midrule
v2 frozen   & heads only                          & 145 & \TODO{} & \TODO{} & \TODO{} & \TODO{} \\
v3 partial  & +unfreeze last 4 blocks, +helper mask & 73.6 & \TODO{} & \TODO{} & \TODO{} & \TODO{} \\
v4 boosted  & +neck/hand loss $\times4$            & 73.6 & 56.5 & 73.2 & \TODO{} & \TODO{} \\
v5 shape    & +$\alpha_{\text{shape}}{=}20$ (teacher) & \textbf{74.3} & \textbf{57.5} & \textbf{74.4} & \textbf{54.1} & 90 \\
\bottomrule
\end{tabular}
\caption{Held-out Anny-One results (mm) across the transfer schedule. The v5
row is the standalone test-set evaluation (\texttt{run\_eval\_test.sh},
\texttt{--test\_anny\_n 100}, batch size 1; precision 81\,\%, F1 85\,\%); the
training-time evaluation of the same checkpoint reported 73.6 / 56.5 / 73.2,
the 0.7\,mm difference coming from the evaluation batch size. Earlier rows
are training-time values. \TODO{re-evaluate the v2--v4 checkpoints with the
same script for the open cells.}}
\label{tab:main}
\end{table}

\section{Knowledge distillation results}
\label{sec:kd_results}

All students use a ViT-S backbone, $672{\times}672$ input,
$\lambda_{\text{kd}}{=}1$, $T{=}4$, channel-wise softmax, and the same
\texttt{anny\_s2\_shape\_v5} teacher. Table~\ref{tab:student_runs} lists the
configurations; each adds one change to the previous one so that effects are
attributable.

\paragraph{Feature-only distillation (v1/v2).} The first student had randomly
initialised heads and an ImageNet-only DINOv2-S backbone. Run v1 (300 epochs,
LR $10^{-4}$ halved every 30 epochs) reduced $\mathcal{L}_{\text{kd}}$ from 2679
to $\sim$46 but stopped improving after epoch $\sim$60: the schedule had
decayed the LR to $3.9{\cdot}10^{-7}$ by epoch 240. Run v2 resumed for 200
epochs at $2{\cdot}10^{-5}$ halved every 60 and the KD loss kept falling
(463$\rightarrow$38), but the student converged to \textbf{193.4\,mm PVE}
against the teacher's 74.3\,mm --- 2.6$\times$ worse. The Multi-HMR authors'
own ViT-S trails their ViT-L by only 9\,\% at the same resolution, so this gap
is a defect of the distillation setup, not of model capacity: the KL term
supervises only the backbone, and the heads, which turn features into a body,
received nothing from the teacher.

\begin{table}[h]
\centering\small
\begin{tabular}{llll}
\toprule
Run & Student backbone init & Heads & Added \\
\midrule
v1/v2 & ImageNet DINOv2-S & random & feature KD \\
v3 & ImageNet DINOv2-S & 54/67 from teacher & output KD \\
v4 & \texttt{multiHMR\_672\_S} & 54/67 from teacher & colour jitter, batch 8 \\
v5 & \texttt{multiHMR\_672\_S} & 54/67, frozen 5 epochs & 2D losses from epoch 5 \\
v6 & \texttt{multiHMR\_672\_S} & 67/67 via projector & \texttt{--head\_dim 1024} \\
v6b & continues v6 & & 150 more epochs, LR $2{\cdot}10^{-5}$ \\
\bottomrule
\end{tabular}
\caption{Student configurations. Each row adds to the previous one.}
\label{tab:student_runs}
\end{table}

\paragraph{Head transfer (v3--v6).} Table~\ref{tab:kd_epochs} gives the
held-out PVE and PA-PVE of the four head-transfer runs at equal epochs. Two
things are visible. First, every head-transfer run starts \emph{worse} than
the random-head baseline (377, 298 and 378\,mm against 228\,mm at epoch 0):
trained head blocks reading through random entry projections produce
confident wrong poses. Second, the ordering reverses after roughly 20 epochs
and is then monotone at every evaluation: ImageNet backbone $<$ HMR-trained
backbone $<$ head freeze $<$ projector-fed heads. The HMR-trained student
backbone (v4) is the single largest gain, 36--79\,mm over v3 at equal epochs;
the freeze (v5) helps PA-PVE by 4--9\,mm; and feeding the heads teacher-width
tokens (v6) improves both metrics further and crosses the random-head
baseline's \emph{final} PVE by epoch 60. The pre-registered success test for
v5 (``below 228\,mm by epoch 5--10'') was not met by any run; it was the wrong
test, because the projections need about 20 epochs to adapt.

\begin{table}[h]
\centering\small
\begin{tabular}{lrrrrr}
\toprule
Epoch & v3 & v4 & v5 & v6 & v2 (random heads, final) \\
\midrule
0   & 377 / 190 & 298 / 199 & 378 / 192 & 327 / 172 & \\
30  & 305 / 160 & 239 / 163 & 240 / 155 & \textbf{221 / 143} & \\
60  & 228 / 147 & 197 / 138 & 205 / 136 & \textbf{186 / 127} & 193 / 128 \\
100 & 213 / 136 & 179 / 128 & 175 / 119 & \textbf{175 / 116} & \\
130 & 199 / 134 & 163 / 122 & 159 / 113 & \textbf{155 / 106} & \\
best & 191 / 133 & 156 / 117 & 159 / 113 & \textbf{152 / 105} & \\
\bottomrule
\end{tabular}
\caption{Held-out PVE / PA-PVE (mm) of the head-transfer students at equal
epochs. Teacher: 74.3 / 57.5.}
\label{tab:kd_epochs}
\end{table}

\paragraph{Continuation (v6b).} v6 reached its 48-hour limit at epoch 138
with the loss still falling; 138 epochs of 8000 images are under two passes
over Anny-One. Continued for 150 epochs at $2{\cdot}10^{-5}$ (with the
throughput changes of Section~\ref{sec:kd}, 20 hours instead of 48), the
student reached its best value at epoch 130, \textbf{126.9\,mm PVE / 92.1\,mm
PA-PVE} (MPJPE 132.0), after which the curve flattened within $\pm4$\,mm at an
LR of $5{\cdot}10^{-6}$ (Table~\ref{tab:kd_eval}). Detection precision rose
from 71 to 78\,\% over the same run (teacher: 81\,\%); recall stayed at
87--89\,\% (teacher: 90\,\%).

\paragraph{Batch size (v6c).} The same continuation at batch 16 on an 80\,GB
GPU, i.e.\ twice the images per epoch, reached \textbf{114.5\,mm / 84.6\,mm}
after 150 epochs and was still improving at the end ($-7$\,mm over its last
five epochs), 12\,mm better than v6b. The student was never
capacity-limited; it was limited by the number of images it had seen. The
cost is a higher NaN-gradient skip rate (0.7\,\% of steps against 0.4\,\%).
The best synthetic-only student therefore stands at 1.5$\times$ the
teacher's error, down from 2.6$\times$ for feature-only distillation.

\begin{table}[h]
\centering\small
\begin{tabular}{llrrrrrr}
\toprule
Model & Backbone & Params & PVE & PA-PVE & MPJPE & Prec. & Recall \\
\midrule
Teacher (\texttt{anny\_s2\_shape\_v5}) & ViT-L & 318.4\,M & 74.3 & 57.5 & 74.4 & 81 & 90 \\
Student v2 (feature KD only) & ViT-S & 32.8\,M & 193.4 & 128.2 & 210.7 & & 86 \\
Teacher v9 (v5 + 30\,\% real 3DPW, ep.\ 62) & ViT-L & 318.4\,M & 78.5 & 57.4 & 78.6 & 83 & 89 \\
Fresh model v10 (new heads, 30\,\% real from start, ep.\ 299) & ViT-L & 318.4\,M & 87.4 & 71.1 & 99.2 & 81 & 89 \\
Student v6b (projector-fed heads, ep.\ 130) & ViT-S & 37\,M & 126.9 & 92.1 & 132.0 & 77 & 88 \\
Student v6c (v6b recipe, batch 16, ep.\ 149) & ViT-S & 37\,M & \textbf{114.5} & \textbf{84.6} & \textbf{118.4} & 81 & 89 \\
Student v7 (v6b + 30\,\% real 3DPW, ep.\ 130) & ViT-S & 37\,M & 120.6 & 88.2 & 125.2 & 78 & 87 \\
Student v8 (v6c + 30\,\% real, teacher v9, ep.\ 52) & ViT-S & 37\,M & 120.6 & 86.9 & 125.0 & 76 & 87 \\
\bottomrule
\end{tabular}
\caption{Student vs.\ teacher on the 100-image Anny-One holdout (mm), all
rows from the standalone evaluation script except v6b (training-time). The
teacher's precision (81--83\,\%) exceeds every student's (76--81\,\%) at
similar recall. A ViT-B student was not trained.}
\label{tab:kd_eval}
\end{table}

\section{SMPL-to-Anny conversion results}
\label{sec:fit_results}

\paragraph{The 34\,mm floor and its cause.} The first converter was validated
on four frames of \texttt{downtown\_bar\_00} (mean 34.5\,mm; stage-wise errors
in Table~\ref{tab:fit}) and then run on four diverse test scenes with 40
frames each. Two facts in that run's log identified a defect before any
render was examined: the error was 30--37\,mm on \emph{every} frame of every
scene, walking or fencing, and all four actors came out with the same body
(female, age 0.67, height 0.33, Asian 0.73). A residual that does not depend
on the pose and a shape that does not depend on the person mean the shape
was absorbing a fixed skeleton-definition mismatch. The rest-pose comparison
of Section~\ref{sec:smpl2anny} then showed the mismatch: \texttt{spine1} mapped
25\,cm too high and the hips 4\,cm off. (The accompanying renders had looked
wrong for an unrelated reason: the renderer had used 3DPW's
\texttt{img\_frame\_ids}, the index map into the 60\,Hz pose arrays, as image
numbers, drawing pose frame $t$ on image $2t$.)

\begin{table}[h]
\centering\small
\begin{tabular}{lrrrr}
\toprule
Frame & root only & +pose & +shape & final (mm) \\
\midrule
0   & 137.9 & 40.4 & 36.2 & 36.2 \\
50  & 131.2 & 43.1 & 35.1 & 35.1 \\
100 & 169.5 & 39.4 & 34.1 & 34.1 \\
199 & 135.0 & 39.0 & 32.6 & 32.6 \\
\midrule
mean & & & & 34.5 \\
\bottomrule
\end{tabular}
\caption{First converter (name-based correspondence): mean joint error after
each optimisation stage on \texttt{downtown\_bar\_00}.}
\label{tab:fit}
\end{table}

\paragraph{Effect of each correction.} Table~\ref{tab:fit_versions} isolates
the two corrections on the hardest test scene, \texttt{outdoors\_fencing\_01}
(38 frames at stride 25, including full lunges). Correcting the
correspondence and the targets (v2) removed the common floor: limbs and head
dropped to 1--10\,mm, but the torso joints stayed at 43--73\,mm because of the
definition gap, and the optimiser compromised by displacing the pelvis 5\,cm
and driving \texttt{proportions} to 0.06. The virtual joints (v2.1) brought
the torso to 13--22\,mm, the pelvis to 19\,mm and \texttt{proportions} to 0.27,
leaving the arm chain (shoulders 23--36\,mm, wrists 25--27\,mm) as the largest
remaining term. The per-joint table is flat, which is the signature of a
correct correspondence.

\begin{table}[h]
\centering\small
\begin{tabular}{p{5.6cm}rrp{4.6cm}}
\toprule
Converter & mean (mm) & max & worst joints (mm) \\
\midrule
v1: name-based mapping, neutral SMPL targets & 34.3 & 37 & all joints, common floor \\
v2: geometry-checked mapping, 3DPW joints, fixed gender & 24.6 & 30 & spine1 73, collars 55--59, spine3 58, pelvis 53 \\
v2.1: + virtual torso joints, bone mask, warm start & \textbf{15.8} & 18.8 & shoulder 36, wrists 25--27 \\
\bottomrule
\end{tabular}
\caption{Mean joint error over all 24 joints on \texttt{outdoors\_fencing\_01}
(38 frames) for the three converter versions.}
\label{tab:fit_versions}
\end{table}

\paragraph{Diverse scenes.} With v2.1, the four test scenes chosen for
variety convert to 15.8--20.1\,mm (Table~\ref{tab:fit_scenes}), with no frame
near the 80\,mm rejection threshold except one occluded office frame
(43.5\,mm), and the fitted bodies now differ between actors. One number did
not: \texttt{age} came out at exactly 0.67 for all four, as it had in every
earlier version. Joint positions cannot distinguish a 25-year-old from a
45-year-old; the optimiser was using the age morph to absorb a residual
torso-length gap. Age is therefore fixed at 0.55 in the final converter, as
described in Section~\ref{sec:smpl2anny}.

\begin{table}[h]
\centering\small
\begin{tabular}{lrrrrrr}
\toprule
Scene & frames & mean (mm) & max & height & proportions & age \\
\midrule
\texttt{downtown\_walking\_00} & 40 & 17.5 & 19.0 & 0.52 & 0.34 & 0.67 \\
\texttt{outdoors\_fencing\_01} & 38 & 15.8 & 18.8 & 0.51 & 0.27 & 0.67 \\
\texttt{office\_phoneCall\_00} & 32 & 20.1 & 43.5 & 0.49 & 0.42 & 0.67 \\
\texttt{flat\_guitar\_01} & 30 & 16.2 & 18.0 & 0.52 & 0.23 & 0.67 \\
\bottomrule
\end{tabular}
\caption{Converter v2.1 on four diverse 3DPW test scenes (stride 25, shared
body per scene, all 24 joints). Age was still a free phenotype in this run.}
\label{tab:fit_scenes}
\end{table}

\paragraph{Whole dataset.} The final converter (v2.2: age fixed, mesh
bounding boxes) was run on all 87 person--sequences at every third frame.
On the training-split sequences inspected so far the mean error is
12.5--22.9\,mm per sequence (e.g.\ \texttt{courtyard\_arguing\_00}: 255 frames,
12.8\,mm mean, 15.3\,mm max), the hardest single frames are acrobatic
(capoeira, 70.8\,mm) and remain below the rejection threshold, and no fit
failed. The two recurring courtyard actors receive identical bodies across
their sequences (male: height 0.60, proportions 0.41; female: 0.34, 0.70),
whereas the distinct test-split actors receive distinct ones, which is the
behaviour one would want from a shape estimate. The merged training labels
comprise 4\,451 images of the train split. Table~\ref{tab:labels} gives the
full count: 16\,631 labelled images with 20\,028 people across the three
splits, every one of them under the 80\,mm rejection threshold; the only
drops are 162 validation-split people who are mostly outside the frame.

\begin{table}[h]
\centering\small
\begin{tabular}{lrrrl}
\toprule
Split & sequences & images & people & dropped \\
\midrule
train & 24 & 4\,451 & 5\,183 & 0 \\
validation & 12 & 2\,458 & 2\,995 & 162 (visible $<$ 0.3) \\
test & 24 & 9\,722 & 11\,850 & 0 \\
\midrule
total & 60 & 16\,631 & 20\,028 & 162 \\
\bottomrule
\end{tabular}
\caption{Converted 3DPW labels (every third frame, converter v2.2). No frame
exceeded the 80\,mm fit-error threshold.}
\label{tab:labels}
\end{table}

Because the ground truth itself carries this residual, any Anny-space
accuracy measured on 3DPW is bounded below by roughly 13--20\,mm; this is
stated wherever such numbers are reported.

\section{Qualitative results}

\TODO{Figure: in-domain Anny-One holdout images with predicted meshes
(current best checkpoint). Figure: real photographs (tennis, stage, party
images used throughout development). Figure: 3DPW frames. Note visibly the
neck elongation on all, and correct body pose after the weight-loading fix.}

\TODO{Figure: converted 3DPW ground truth rendered over the original frames
(\texttt{threedpw\_anny\_v21/renders/}): the Anny mesh with SMPL joints (red)
and fitted virtual joints (green), e.g.\ the fencing lunge and the seated
guitar frame, which are the supervisor's qualitative acceptance check.}

% =============================================================================
\chapter{Ablation Studies and Root-Cause Investigations}
\label{ch:ablation}
% =============================================================================

The supervisor's guideline for this chapter is to justify each design choice by
removing it under otherwise identical settings and reporting the change. Two of
the investigations below go further: they isolate the cause of an observed
defect by testing and eliminating alternative hypotheses.

\section{Shape loss weight}
\label{sec:shape_ablation}

\paragraph{Observation.} The trained model predicted nearly the same phenotypes
for every person. Comparing predictions against Anny-One ground truth on the
same held-out samples (Table~\ref{tab:shape}) showed three phenotypes offset by
4--6 standard deviations of the data: the shape head was regressing dataset
averages, not per-person values.

\paragraph{Hypothesis.} With $\alpha_{\text{shape}}{=}1$ against
$\alpha_{\text{j3d}}{=}\alpha_{\text{v3d}}{=}100$, the network gains far more by
adjusting geometry than by getting shape right, so the shape head is
under-trained.

\paragraph{Test.} Raise $\alpha_{\text{shape}}$ to 20 and continue training
(run v5), all else identical.

\begin{table}[h]
\centering\small
\begin{tabular}{lrrrrr}
\toprule
Phenotype & GT mean & GT std & v4 error ($\alpha{=}1$) & v5 error ($\alpha{=}20$) & \\
\midrule
age          & 0.553 & 0.036 & $+0.199$ & $\mathbf{-0.002}$ & fixed \\
height       & 0.579 & 0.026 & $-0.154$ & $\mathbf{-0.026}$ & fixed \\
proportions  & 0.763 & 0.070 & $-0.017$ & $\mathbf{-0.001}$ & fixed \\
firmness     & 0.757 & 0.039 & $-0.167$ & $-0.190$ & unchanged \\
gender       & 0.694 & --    & $-0.042$ & $-0.063$ & $\approx$ \\
weight       & 0.382 & --    & $+0.056$ & $+0.064$ & $\approx$ \\
\bottomrule
\end{tabular}
\caption{Predicted-minus-ground-truth phenotype error on the same Anny-One
holdout samples, before and after raising the shape loss weight from 1 to 20.
Errors are in phenotype units ($[0,1]$).}
\label{tab:shape}
\end{table}

\paragraph{Result.} Age error fell by 99\,\% and height error by 83\,\%;
\texttt{proportions} became essentially exact. \texttt{firmness} did not
respond, and is discussed in Chapter~\ref{ch:discussion}. A sweep of
$\alpha_{\text{shape}}\in\{1,5,10,20,50\}$ under identical settings would
locate the optimum; \TODO{run and tabulate PVE and shape error per value.}

\section{Neck elongation: root-cause isolation}
\label{sec:neck}

Every rendered mesh showed a visibly elongated, thin neck. Four candidate
causes were tested (Table~\ref{tab:neck}); each test was designed so that its
outcome would unambiguously confirm or eliminate the hypothesis.

\begin{table}[h]
\centering\small
\begin{tabular}{p{3.2cm}p{5.2cm}p{4.6cm}}
\toprule
Hypothesis & Test & Result \\
\midrule
The Anny body model itself has a long neck &
Render every person in rest pose (all rotations identity, predicted shape kept) &
Neck normal $\rightarrow$ eliminated \\
Helper-bone mask removes neck bones from the loss, so the $\times4$ boost multiplied zero &
Script listing masked/active status of every neck bone &
\texttt{neck01/02/03}, \texttt{head} all active $\rightarrow$ eliminated \\
Domain gap (synthetic $\rightarrow$ real photos) &
Run on Anny-One's own held-out images &
Elongation present in-domain $\rightarrow$ eliminated \\
Shape head over-predicts \texttt{proportions} &
GT vs.\ predicted phenotypes on the same samples &
\texttt{proportions} accurate to 0.018 (v4) / 0.001 (v5) $\rightarrow$ eliminated \\
\bottomrule
\end{tabular}
\caption{Elimination of candidate causes for the neck artifact.}
\label{tab:neck}
\end{table}

The surviving explanation is linear blend skinning across Anny's three stacked
neck bones plus head: when the network predicts slightly different rotations
for each, the skin between them stretches. This is consistent with every
observation: rest pose is normal (no rotation differential), the $\times4$ neck
loss boost did not change it (rotations are already near the target), it appears
in-domain, and shape is accurate. It is a property of the rig, not a training
failure. A side finding is that an inference-time ``fix'' forcing
\texttt{proportions} to 0.5 was in fact pushing it 0.26 away from the true
dataset mean of 0.763 and was removed.

\section{Per-joint loss boosting}

Up-weighting neck and hand bones by $\times4$ in
$\mathcal{L}_{\text{rot}}$ and $\mathcal{L}_{\text{j3d}}$ (run v4, from the v3
checkpoint, identical settings otherwise) left holdout PVE unchanged at
73.6\,mm (best 72.0 at epoch 125) and did not alter the neck. The mask check
above confirms the boost did reach the intended bones; it simply had nothing to
correct. The boost is therefore disabled in the final configuration.

\section{Learning-rate schedule for distillation}

Two otherwise identical distillation runs differ only in schedule. With a halving
every 30 epochs over 300 epochs the learning rate reaches $3.9{\cdot}10^{-7}$ by
epoch 240 and all losses plateau from epoch $\sim$60; with a halving every 60
epochs over 200 epochs the KD loss falls monotonically to the end
(Table~\ref{tab:kd_epochs}). The schedule must be sized to the epoch budget.

\section{Backbone transfer and fine-tuning depth}

\begin{table}[h]
\centering\small
\begin{tabular}{lrr}
\toprule
Configuration & Trainable params & Holdout PVE \\
\midrule
Frozen backbone, heads only & $\sim$14\,M & 145 \\
Last 4 of 24 blocks unfrozen (+helper mask) & $\sim$64\,M & 73.6 \\
Fully unfrozen (from scratch, all fixes) & 318\,M & \TODO{run v7} \\
\bottomrule
\end{tabular}
\caption{Effect of how much of the backbone is trained. The v2$\rightarrow$v3
step also introduced helper-bone masking; a clean separation requires re-running
v3 without the mask. \TODO{}}
\label{tab:ft}
\end{table}

\section{Distilling the heads, not only the backbone}
\label{sec:head_ablation}

\paragraph{Observation.} With feature-level distillation only, the student's
backbone tokens converge toward the teacher's ($\mathcal{L}_{\text{kd}}$
2679$\rightarrow$38) yet PVE plateaus at 193\,mm. The backbone learned the
teacher's representation; the randomly initialised heads did not learn to use
it as well as the teacher's heads do.

\paragraph{Intervention and its side effect.} Copying the 54 shape-compatible
head tensors from the teacher (v3) \emph{raised} first-eval PVE from 228 to
377\,mm. The 13 non-transferable tensors are the head's input paths, so this
is the expected behaviour of a trained function applied to random inputs:
confident, wrong outputs, whose gradients then erode the transferred weights.

\paragraph{Three hypotheses, three runs.} Each was tested by one change
(Table~\ref{tab:kd_epochs}). (i) \emph{The gradients damage the heads}: the
five-epoch freeze (v5) kept the rotation loss at 90 where v4's sat at 300, so
the damage is real, yet PA-PVE at the unfreeze evaluation did not move
(192$\rightarrow$193); protecting the heads is necessary but not the main
effect. (ii) \emph{The backbone, not the heads, is the problem}: initialising
the student backbone from the authors' HMR-trained ViT-S instead of ImageNet
(v4) gave the largest single improvement at every epoch. (iii) \emph{The heads
are fed the wrong input}: moving the KD projector in front of the heads so
that all 67 tensors transfer (v6) gave the best run at every epoch from 10
onward and the best final student. The sequence shows the right diagnosis
sitting one step behind each intervention: v3 fixed the head initialisation
and exposed the entry projections; v5 fixed the damage and exposed that the
projections themselves were the bottleneck; v6 removed them.

\paragraph{What remains.} At 114.5\,mm the best synthetic-only student is
1.5$\times$ the teacher, against a 9\,\% gap between the authors' own ViT-S
and ViT-L. The batch-16 continuation (v6c) took 12\,mm off v6b and was still
falling when its budget ended, so under-training rather than capacity is the
better-supported explanation for most of the remaining gap on synthetic data.
On real photographs the picture is different: there the student plateaus at
the same value whichever teacher distils it and whichever checkpoint it starts
from (Section~\ref{sec:mixed}), i.e.\ the ViT-S itself is the limit.

\section{Student backbone size and KD hyper-parameters}

A ViT-B student and sweeps of $\lambda_{\text{kd}}$ and $T$ were planned but
not run within the time available; all students use $\lambda_{\text{kd}}{=}1$,
$T{=}4$. Given that the ViT-S student is limited by training length on
synthetic data and by its own capacity on real data, a ViT-B student is the
natural next experiment and the recipe transfers unchanged
(\texttt{--backbone dinov2\_vitb14}, \texttt{--head\_dim 1024}).

\section{Synthetic-only vs.\ mixed synthetic--real training}
\label{sec:mixed}

This is the result that addresses the domain-shift concern raised by the
supervisor: a student can match the teacher on synthetic renders while both
rely on render-specific cues absent from photographs. Four models were trained
with 30\,\% of every batch drawn from the 4\,451 converted 3DPW-train images:
the teacher continued from v5 (\textbf{v9}: last four ViT blocks unfrozen, LR
$10^{-5}$), the student continued from v6b with teacher v5 (\textbf{v7}), the
student continued from v6c with teacher v9 (\textbf{v8}, the final student),
and a fresh model with heads trained from scratch on the Naver backbone with
real data present from the start (\textbf{v10}, 300 epochs). All were
evaluated every five epochs on the Anny-One holdout and on 300 converted
3DPW-validation images; model selection used the validation split only.

\paragraph{Validation trajectories.} The real data acts within the first
five epochs. The teacher's 3DPW-validation PVE fell from 127.1 to 78.0\,mm at
epoch 5 and never improved afterwards (80--83\,mm by epoch 150, a sign of
over-fitting 4\,451 images); the student's fell from 149.0 to 116.8 at epoch 5
and settled at 106--108 from epoch 20 onward. On the synthetic holdout the
teacher paid 3--9\,mm, whereas the student first lost 8\,mm and then recovered
to below its starting value (128.5 $\rightarrow$ 120.1), so at $\rho{=}0.3$ the
real data was close to free for the student. v8 reproduced v7 almost exactly
on real photographs (105.7 vs 106.0 at best) despite starting from a 12\,mm
better checkpoint and distilling from a teacher that is 49\,mm better on real
data: the student's real-photo error does not depend on the teacher's domain
gap, which locates the limit in the ViT-S itself. v10 plateaued from epoch
$\sim$100 at 87\,mm synthetic / 90\,mm real on validation, short of the
continued teacher on both (discussed with the test results below).

\paragraph{Test split.} Table~\ref{tab:test3dpw} gives the one-time
evaluation on all 9\,722 converted 3DPW-test images. Every model that saw
real data improved on real photographs: the teacher by 14\,\% (119.1
$\rightarrow$ 102.5\,mm, PA-PVE 68.8 $\rightarrow$ 62.7) and the student by
11\,\% (139.5 $\rightarrow$ 123.6, PA-PVE 88.0 $\rightarrow$ 81.2), for a cost
of 4 and 6\,mm on the synthetic holdout. The gains are, however, about half of
what the validation split had shown (38\,\% and 29\,\%), and the mixed
student no longer beats the synthetic-only teacher on real photographs (123.6
vs 119.1\,mm) as it did on validation (106 vs 127). The reason is in the data:
3DPW's train and validation splits were filmed with the same actors at the
same courtyard locations, whereas the test split contains different people in
different (downtown) scenes. The validation numbers therefore partly measured
how well the models learned those actors; the test numbers measure
generalisation, and are the ones to report. Holding the test split back from
every decision is what made this visible.

\paragraph{Real data from the start (v10).} The fresh model trained with real
data from its first epoch reaches 111.5\,mm on 3DPW-test, 8\,mm better than
the synthetic-only teacher but 9\,mm worse than the teacher that received the
same real data at the end of its staged training (v9), and 87.4\,mm on the
synthetic holdout, 13\,mm worse than either. Its real-photo PA-PVE (76.2) is
also worse than the synthetic-only teacher's (68.8): what the real data gave
it is better global placement and scale, not better body pose and shape. Its
recall on 3DPW (85\,\% against 95--100\,\%) and the unusually large gap
between its joint and vertex errors (MPJPE 99.2 vs PVE 87.4 on synthetic data,
where the other models are within 1--5\,mm) both point the same way: a single
300-epoch run does not replace the staged teacher training, and real data is
best added as a final fine-tuning step.

\begin{table}[h]
\centering\small
\begin{tabular}{llrrrrr}
\toprule
Model & Backbone & real data & PVE & PA-PVE & MPJPE & PA-MPJPE \\
\midrule
Teacher v5 & ViT-L & -- & 119.1 & 68.8 & 124.1 & 63.1 \\
Teacher v9 (ep.\ 62) & ViT-L & 30\,\% & \textbf{102.5} & \textbf{62.7} & \textbf{102.1} & \textbf{56.1} \\
Fresh model v10 & ViT-L & 30\,\% from start & 111.5 & 76.2 & 121.8 & 84.8 \\
Student v6c & ViT-S & -- & 139.5 & 88.0 & 144.1 & 83.1 \\
Student v7 & ViT-S & 30\,\% & 124.5 & 81.7 & 126.2 & 76.8 \\
Student v8 (final) & ViT-S & 30\,\% & \textbf{123.6} & \textbf{81.2} & 125.6 & \textbf{76.3} \\
\bottomrule
\end{tabular}
\caption{3DPW-test (9\,722 converted images, Anny-space metrics in mm). The
13--20\,mm conversion residual is a floor on every row. Recall is 95--100\,\%
for all models except v10 (85\,\%); precision (37--59\,\%) is not meaningful here because 3DPW
annotates only one or two actors per frame and the downtown scenes contain
unannotated bystanders, so each correctly detected passer-by counts as a false
positive.}
\label{tab:test3dpw}
\end{table}

\paragraph{Reading the final model.} The final student (v8: 37\,M
parameters, 8.6$\times$ smaller than the teacher) scores 123.6\,mm on real
photographs against 102.5 for the mixed teacher and 119.1 for the
synthetic-only teacher, and 120.6\,mm on synthetic data against 74.3 / 78.5.
Against the Multi-HMR authors' own 9\,\% ViT-S/ViT-L gap on SMPL-X, the
remaining gap here (21\,mm on real, 42\,mm on synthetic) says that the
distillation recipe is not yet at parity; the ablations above locate the
synthetic-side gap in training length and the real-side gap in the student's
capacity.

\section{Joint correspondence and virtual joints in the converter}
\label{sec:fit_ablation}

Table~\ref{tab:fit_versions} is itself an ablation: the correspondence fix
accounts for 10\,mm of the improvement and the virtual joints for a further
9\,mm, and the two act on different joints (the first on limbs and head, the
second on the torso). Two further observations qualify as negative results.
Down-weighting the torso joints and reporting only a ``reliable'' subset,
tried in an intermediate version, hid the symptom without touching the cause
and should not be used. And the \texttt{age} phenotype, free in every version
up to v2.1, converged to 0.67 for every actor in every scene; it is not
identifiable from joints and is now fixed. Each converter flag
(\texttt{--offset\_joints}, \texttt{--free\_bones}, \texttt{--shape\_reg},
\texttt{--warm\_start}) can be switched off individually to reproduce the
intermediate rows.

% =============================================================================
\chapter{Discussion}
\label{ch:discussion}
% =============================================================================

\paragraph{What improved.} The staged transfer schedule halved the held-out
PVE (145$\rightarrow$74.3\,mm) and the shape re-weighting corrected the
systematic phenotype bias almost entirely for age, height and proportions. The
distillation framework works end-to-end: an 8.6$\times$ smaller student
reaches 127\,mm PVE, down from 193\,mm with feature-only distillation, once
the teacher's heads are transferred through a projector that feeds them
teacher-width tokens. The SMPL-to-Anny converter went from a 34\,mm floor that
was the same for every frame and every person to 13--20\,mm with bodies that
differ between actors, and the whole of 3DPW is now available as Anny-space
labels.

\paragraph{What did not improve, and best guesses why.}
\begin{itemize}
  \item \emph{Neck elongation.} Unchanged by masking, loss boosting, shape
        correction or in-domain evaluation. The best explanation is LBS stretch
        across the stacked neck bones. Because it is a rig property, the
        remedies lie outside training: fewer neck bones, adjusted skinning
        weights, or a smoothness constraint across the neck chain in the loss.
  \item \emph{\texttt{firmness}.} Did not respond to a 20$\times$ increase in
        shape supervision and does not move under joint-based fitting either.
        Best guess: it is a soft-tissue property with almost no visual
        signature at $672$ pixels on clothed people, and none in joint
        positions, i.e.\ it is unobservable from the available signal rather
        than under-trained.
  \item \emph{PVE plateau at $\sim$73\,mm.} Every continuation converges to the
        same value. Best guess: the frozen early backbone blocks limit
        adaptation; the fully-unfrozen v7 run (\TODO{}) tests this.
  \item \emph{The remaining student--teacher gap (127 vs 74.3\,mm).}
        Matching backbone features was necessary but not sufficient; the
        heads had to be transferred, and transferring them across a width
        mismatch introduced its own failure (random entry projections feeding
        trained blocks) that a warm-up alone did not cure and the
        projector-fed design did. Each step was correct about the defect it
        fixed and wrong about the next one it exposed
        (Tables~\ref{tab:student_runs} and \ref{tab:kd_epochs}). The batch-16
        continuation showed that the synthetic-side gap is mostly training
        length; the mixed runs showed that the real-side gap is not (two
        different teachers and starting points gave the same real-photo
        error), so it belongs to the student's capacity.
\end{itemize}

\paragraph{Pipeline bugs whose symptoms mimic model failures.}
\label{sec:pipeline}
Several weeks were spent on artifacts that turned out to be engineering
defects. They are documented because their symptoms are easily misread as
learning failures:
\begin{itemize}
  \item The demo selected the model class by checkpoint \emph{filename};
        Anny checkpoints were routed to a class with different module names,
        and \texttt{strict=False} loading silently loaded zero weights, so every
        person rendered in the same mean pose. Fix: route by checkpoint
        contents and assert $\geq$80\,\% of tensors loaded.
  \item Two code paths produced translations in centimetres and metres
        respectively; an unconditional $/100$ discarded every in-domain
        detection in one place and placed people 2.8\,cm from the camera in
        another. Fix: unit auto-detection.
  \item A single global field of view was applied to Anny-One images whose
        true FOV spans 51--110$^\circ$, mis-scaling meshes in depth. Fix:
        per-image intrinsics.
  \item Reusing a run name let checkpoint cleanup, which keeps the ten
        highest-numbered files, delete a new run's checkpoints.
  \item A degenerate 6D rotation initialisation (indices 0 and 4 instead of 0
        and 3) produced NaN for every bone before optimisation began in the
        SMPL-to-Anny fit. Two earlier hypotheses (invalid 3DPW frames, the
        SMPL forward pass) were each ruled out by the log before the real
        cause was found.
  \item The converter's overlay renderer used 3DPW's \texttt{img\_frame\_ids}
        as image numbers. That array is the index map into the 60\,Hz pose
        arrays, roughly $2t$, so every mesh was drawn on the frame where the
        person had already walked elsewhere, and the fits looked far worse
        than they were. The tell-tale was in the log (a request for image
        1401 of a 1387-frame sequence), not in the picture.
  \item A joint correspondence chosen by name was wrong by 25\,cm for one
        spine joint because the two rigs number the spine in opposite
        directions. It produced a plausible-looking 34\,mm fit and a
        ``reliable-joint'' metric that quietly excluded the affected joints;
        only comparing the two rest skeletons geometrically exposed it.
  \item A learning-rate schedule copied from short fine-tuning runs (halving
        every 30 epochs) silently froze a 300-epoch distillation run from
        epoch $\sim$60 onward; the loss curves looked converged when the
        optimiser had simply run out of step size.
  \item The 3DPW archive lives on a host path (\texttt{/ds-av}) that was not
        in the job container's mount list; the failure surfaced as an
        unrelated ``file not found'' inside the container.
\end{itemize}
The common lesson is that silent failure is the dominant failure mode in this
class of pipeline; loud, verified failure (load percentages, unit checks, NaN
guards) would have saved most of the time lost.

\paragraph{Limitations.} The in-domain test set is only 100 samples; a
larger holdout requires retraining with a larger \texttt{--val\_anny\_n}. 3DPW
results in Anny space depend on the conversion residual (13--20\,mm), which
bounds the achievable measured accuracy on real data; the arm chain
(shoulders, wrists) is the largest remaining term. The converted labels fix
age and gender from annotations and cannot recover soft-tissue phenotypes,
so the real data supervises shape only through height and proportions. The
3DPW validation split shares actors and locations with the training split,
so it over-states the benefit of real data by about a factor of two relative
to the test split; only the test numbers should be quoted. A ViT-B student
was not trained. The neck artifact is unresolved at the rig level.

% =============================================================================
\chapter{Conclusion}
\label{ch:conclusion}
% =============================================================================

This thesis adapted Multi-HMR to the open Anny body model and trained it on
Anny-One, reaching 74.3\,mm PVE on held-out synthetic data. It identified and
corrected a loss-imbalance defect in shape estimation with quantified
before/after evidence, isolated the cause of a persistent rendering artifact by
eliminating four alternatives, built a knowledge-distillation framework that
compresses the model 8.6$\times$ and, through a projector that lets the
teacher's heads transfer intact, brings the student from 193 to 114.5\,mm PVE
on synthetic data, and produced an SMPL-to-Anny converter whose
correspondence is verified by rest-pose geometry and whose virtual joints
reduce the conversion floor from 34 to 13--20\,mm. With it the whole of 3DPW
was converted into Anny-space labels, and training with 30\,\% real images
reduced the real-photograph error of the teacher by 14\,\% (119 $\rightarrow$
103\,mm on 3DPW-test) and of the student by 11\,\% (140 $\rightarrow$ 124\,mm)
at a cost of 4--6\,mm on synthetic data. Every intervention is exposed as a
flag so the experiments can be re-run under identical settings.

\paragraph{Future work.} Train a ViT-B student with the same recipe, since
the ViT-S is capacity-limited on real photographs; evaluate the final models
under the Multi-HMR SMPL protocol on 3DPW-test for comparison with published
numbers; use the teacher's predictions on unlabelled
real photographs as pseudo-labels (the output-distillation term already
provides the mechanism); extend the virtual-joint treatment to the arm chain;
address the neck at the rig level; and, per the original topic description,
extend the system from images to video with temporal consistency, for which
HMR2.0's tracking formulation is the natural starting point.

% =============================================================================
% Appendix material lives in appendix.tex (hyper-parameter tables per run,
% additional qualitative figures, diagnostic scripts).
% =============================================================================
